\documentclass[10pt,a4paper]{amsart}
\usepackage[margin=19mm]{geometry}
\usepackage{amsmath,amssymb,mathtools}
\usepackage[scr=rsfs,cal=euler]{mathalfa}
\usepackage{xfrac}
\usepackage[unicode,hidelinks,bookmarks=false]{hyperref}
\newcommand{\ie}{i.e.\ }
\newcommand{\cf}{cf.\ }
\newcommand{\resp}{respectively\ }
\newcommand{\Subsection}{\subsection}
\providecommand{\nobreakdash}{\nobreak\hbox{-}}
\setlength{\emergencystretch}{2em}
\multlinegap0pt
\usepackage{xcolor}
\usepackage{amssymb}
\newtheorem{theorem}{Theorem}
\newtheorem{lemma}{Lemma}
\title{Elementary Proofs of Two Congruences for Partitions with Odd Parts Repeated at Most Twice}
\author{James A. Sellers}
\date{Source: arXiv:2409.12321v2 (CC BY 4.0)}
\begin{document}
\maketitle

\noindent\emph{Source and licence.} Elementary Proofs of Two Congruences for Partitions with Odd Parts Repeated at Most Twice. Version arXiv:2409.12321v2 (CC BY 4.0), distributed by arXiv under the Creative Commons Attribution 4.0 International licence, http://creativecommons.org/licenses/by/4.0/, as recorded in the arXiv licence metadata for this exact version and shown on the abstract page https://arxiv.org/abs/2409.12321v2. Full licence notice: \texttt{licenses/arxiv-2409.12321-CC-BY-4.0.txt}.\par\smallskip

\noindent\textit{Editorial scope.} Counted unit: the article's theorem JAS\_id1 (source0.tex lines 148--154). The article's complete original proof of it is delivered with it (source0.tex lines 201--248, one \texttt{proof} environment of 48 lines, which the article places away from the statement). No mathematics is written, repaired or re-derived: the statement and the proof are the article's own lines, in the article's own notation, and no retained line is substituted or re-flowed.  Retained uncounted as the source-local context the proof needs: lines 99--108; lines 156--163; lines 174--180; lines 187--193.  Nothing was omitted from any retained span, so the package carries no omission record.  No retained line contains a citation, so the wrapper carries no bibliography and no bibliography entry.  No retained line contains a figure environment, an image-inclusion command or drawing code, so no image is delivered and the package declares no asset dependency.  Editorial formatting only, in addition: an AMS \texttt{amsart} preamble that restates the article's own declarations and macros used by the retained lines, namely the environments \texttt{theorem}, \texttt{lemma} on separate counters, together with the standard packages those lines need.  The retained environments print in this document as Theorem 1, Lemma 1 in order of appearance, whereas the article prints the same items with the numbering of its own full document.  The complete native source of the article as distributed in the arXiv e-print for this exact version is bundled unchanged as \texttt{sources/165471/source0.tex}.
\begin{equation}
\label{genfn_a}
\sum_{n\geq 0} a(n)q^n 
=
\prod_{i\geq 1} \frac{(1-q^{6i-3})}{(1-q^i)} 
= 
\prod_{i\geq 1} \frac{(1-q^{6i-3})}{(1-q^i)} \cdot \frac{(1-q^{6i})}{(1-q^{6i})} 
= 
\frac{f_3}{f_1f_6}.
\end{equation}

\begin{theorem}
\label{JAS_id2}
We have 
$$
\sum_{n\geq 0} a(4n+3)q^n = 2\frac{f_2^4f_8^2f_{12}}{f_1^5f_4^2f_6}.
$$

\end{theorem}

\begin{lemma}
\label{lemma_f3overf1}
We have 
$$
\frac{f_3}{f_1} = \frac{f_4f_6f_{16}f_{24}^2}{f_2^2f_8f_{12}f_{48}} + q\frac{f_6f_8^2f_{48}}{f_2^2f_{16}f_{24}}.
$$
\end{lemma}

\begin{lemma}
\label{lemma_1overf1squared}
We have 
$$
\frac{1}{f_1^2} = \frac{f_8^5}{f_2^5f_{16}^2} + 2q\frac{f_4^2f_{16}^2}{f_2^5f_8}.
$$
\end{lemma}

\begin{theorem}
\label{JAS_id1}
We have 
$$
\sum_{n\geq 0} a(4n+2)q^n = 2\frac{f_2f_6^2f_8^2}{f_1^4f_3f_{12}}.
$$
\end{theorem}

\begin{proof}(of Theorem \ref{JAS_id1} and Theorem \ref{JAS_id2}) 
We begin by finding the 2--dissection of the generating function for $a(n)$, using (\ref{genfn_a}) as our starting point.  
\begin{align*}
\sum_{n\geq 0} a(n)q^n 
&= 
\frac{f_3}{f_1f_6} \\
&= 
\frac{1}{f_6}\left( \frac{f_3}{f_1} \right) \\
&= 
\frac{1}{f_6}\left( \frac{f_4f_6f_{16}f_{24}^2}{f_2^2f_8f_{12}f_{48}} + q\frac{f_6f_8^2f_{48}}{f_2^2f_{16}f_{24}} \right) \textrm{\ \ \ (Lemma \ref{lemma_f3overf1}) } \\
&= \frac{f_4f_{16}f_{24}^2}{f_2^2f_8f_{12}f_{48}} + q\frac{f_8^2f_{48}}{f_2^2f_{16}f_{24}}.
\end{align*}
Thus, we can immediately read off the following:  
\begin{align*}
\sum_{n\geq 0} a(2n)q^n &= \frac{f_2f_{8}f_{12}^2}{f_1^2f_4f_{6}f_{24}},  \\
\sum_{n\geq 0} a(2n+1)q^n &= \frac{f_4^2f_{24}}{f_1^2f_{8}f_{12}}.
\end{align*}
We now perform two additional 2--dissections.  First, 
\begin{align*}
\sum_{n\geq 0} a(2n)q^n 
&= 
\frac{f_2f_{8}f_{12}^2}{f_4f_{6}f_{24}}\left( \frac{1}{f_1^2} \right)\\
&= 
\frac{f_2f_{8}f_{12}^2}{f_4f_{6}f_{24}}\left( \frac{f_8^5}{f_2^5f_{16}^2} + 2q\frac{f_4^2f_{16}^2}{f_2^5f_8}  \right)  \textrm{\ \ \ (Lemma \ref{lemma_1overf1squared}) } \\
&= 
\frac{f_8^6f_{12}^2}{f_2^4f_4f_6f_{16}^2f_{24}} + 2q\frac{f_4f_{12}^2f_{16}^2}{f_2^4f_6f_{24}}.
\end{align*}
We then see that 
$$
\sum_{n\geq 0} a(4n+2)q^n = 2\frac{f_2f_{6}^2f_{8}^2}{f_1^4f_3f_{12}}
$$
and this proves Theorem \ref{JAS_id1}.  
Finally, 
\begin{align*}
\sum_{n\geq 0} a(2n+1)q^n 
&=
\frac{f_4^2f_{24}}{f_{8}f_{12}} \left( \frac{1}{f_1^2} \right)\\
&= 
\frac{f_4^2f_{24}}{f_{8}f_{12}}\left(  \frac{f_8^5}{f_2^5f_{16}^2} + 2q\frac{f_4^2f_{16}^2}{f_2^5f_8}  \right) \\
&= 
\frac{f_4^2f_8^4f_{24}}{f_2^5f_{12}f_{16}^2} + 2q\frac{f_4^4f_{16}^2f_{24}}{f_2^5f_8^2f_{12}}.
\end{align*}
Hence, 
$$
\sum_{n\geq 0} a(4n+3)q^n = 2\frac{f_2^4f_{8}^2f_{12}}{f_1^5f_4^2f_{6}}
$$
and this proves Theorem \ref{JAS_id2}.  
\end{proof}

\end{document}
